% Luc Destombe --- licence CC BY-NC-SA 4.0
% https://creativecommons.org/licenses/by-nc-sa/4.0/deed.fr
% Fichier autonome : compiler avec pdflatex, deux fois (signets, nombre de pages).
\def\lyceeversion{1.0}
\documentclass[a4paper,10pt]{article}
\makeatletter
\newif\iflycee@niveaux
\lycee@niveauxtrue
\newif\iflycee@cm
\newif\iflycee@palatino

\RequirePackage[utf8]{inputenc}
\RequirePackage[T1]{fontenc}
\RequirePackage[french]{babel}
\RequirePackage{etoolbox}

\newcommand{\ShorthandsOff}{\@ifundefined{shorthandoff}{}{\shorthandoff{;:!?}}}
\newcommand{\ShorthandsOn}{\@ifundefined{shorthandon}{}{\shorthandon{;:!?}}}
\ShorthandsOff
\AtBeginDocument{\ShorthandsOn}
\AtBeginEnvironment{tikzpicture}{\ShorthandsOff}

\RequirePackage{geometry}
\RequirePackage{fancyhdr}
\RequirePackage{lastpage}
\RequirePackage{multicol}
\RequirePackage{array}
\RequirePackage{enumitem}
\RequirePackage{xcolor}
\RequirePackage{needspace}

\RequirePackage{amsmath,amssymb}
\iflycee@cm\else\RequirePackage{fourier}\fi
\DeclareMathAlphabet{\mathcal}{OMS}{cmsy}{m}{n}
\renewcommand{\ttdefault}{lmtt}
\begingroup\catcode`\_=\active \catcode`\^=\active
\global\let_\sb \global\let^\sp
\endgroup
\AtBeginDocument{\catcode`\_=12 \mathcode`\_="8000
  \catcode`\^=12 \mathcode`\^="8000 }
\newlength{\retraitcalculs}
\setlength{\retraitcalculs}{2em}

\RequirePackage{tikz}
\usetikzlibrary{arrows.meta,calc,babel,shapes.misc}
\RequirePackage[most]{tcolorbox}
\tcbuselibrary{skins,breakable}

\definecolor{coulDef}{HTML}{1F4E79}
\definecolor{coulProp}{HTML}{7030A0}
\definecolor{coulRem}{HTML}{C55A11}
\definecolor{coulEx}{HTML}{2E7D32}
\definecolor{coulExo}{HTML}{B03A2E}
\definecolor{coulPart}{HTML}{1F4E79}
\definecolor{coulMeth}{HTML}{1B6E4F}
\definecolor{coulCourbe}{HTML}{0B5ED7}
\definecolor{coulCourbeB}{HTML}{C00000}
\definecolor{coulCourbeC}{HTML}{2E7D32}
\definecolor{coulGrille}{HTML}{8C8C8C}
\definecolor{coulRep}{HTML}{1B6E4F}
\definecolor{coulBar}{HTML}{2E7D32}

\newcommand{\R}{\mathbb{R}}
\newcommand{\Rs}{\mathbb{R}^{*}}
\newcommand{\N}{\mathbb{N}}
\newcommand{\Ns}{\mathbb{N}^{*}}
\newcommand{\Z}{\mathbb{Z}}
\newcommand{\Q}{\mathbb{Q}}
\renewcommand{\le}{\leqslant}
\renewcommand{\ge}{\geqslant}
\renewcommand{\approx}{\simeq}
\newcommand{\vect}[1]{\overrightarrow{#1}}
\newcommand{\norme}[1]{\left\lVert #1 \right\rVert}
\newcommand{\intervalle}[2]{\left[#1\,;#2\right]}
\RequirePackage{cancel}
\renewcommand{\CancelColor}{\color{coulExo}}
\newcommand{\fois}[1]{\times{\color{coulExo}#1}}
\newcommand{\barre}[1]{\cancel{#1}}
\newlength{\lycee@eqh}\newlength{\lycee@eqd}
\newsavebox{\lycee@eqa}\newsavebox{\lycee@eqb}\newsavebox{\lycee@eqc}
\newcommand{\lycee@eqcol}[1]{\hbox{$\renewcommand{\arraystretch}{1.35}%
  \begin{array}[t]{@{}>{\displaystyle}r@{}>{\displaystyle{}}l@{}}#1\end{array}$}}
\newcommand{\lycee@eqmesure}[1]{%
  \ifdim\ht#1>\lycee@eqh\setlength{\lycee@eqh}{\ht#1}\fi
  \ifdim\dp#1>\lycee@eqd\setlength{\lycee@eqd}{\dp#1}\fi}
\newcommand{\lycee@eqtrait}{\hspace{1.2em}\tikz[baseline=0pt]\draw[dash pattern=on 2pt off 2pt,
  line width=0.5pt] (0,\lycee@eqh) -- (0,-\lycee@eqd);\hspace{1.2em}}
\newcommand{\eqproduit}[2]{\par\smallskip\noindent C'est une équation produit :
  $A\times B=0 \iff A=0$ ou $B=0$.\par\nopagebreak\smallskip
  \begingroup\setlength{\lycee@eqh}{0pt}\setlength{\lycee@eqd}{0pt}%
  \setbox\lycee@eqa=\lycee@eqcol{#1}\setbox\lycee@eqb=\lycee@eqcol{#2}%
  \lycee@eqmesure\lycee@eqa\lycee@eqmesure\lycee@eqb%
  \noindent\hspace*{\retraitcalculs}\usebox\lycee@eqa\lycee@eqtrait\usebox\lycee@eqb\endgroup\par\smallskip}
\newcommand{\eqproduittrois}[3]{\par\smallskip\noindent C'est une équation produit :
  $A\times B\times C=0 \iff A=0$ ou $B=0$ ou $C=0$.\par\nopagebreak\smallskip
  \begingroup\setlength{\lycee@eqh}{0pt}\setlength{\lycee@eqd}{0pt}%
  \setbox\lycee@eqa=\lycee@eqcol{#1}\setbox\lycee@eqb=\lycee@eqcol{#2}\setbox\lycee@eqc=\lycee@eqcol{#3}%
  \lycee@eqmesure\lycee@eqa\lycee@eqmesure\lycee@eqb\lycee@eqmesure\lycee@eqc%
  \noindent\hspace*{\retraitcalculs}\usebox\lycee@eqa\lycee@eqtrait\usebox\lycee@eqb\lycee@eqtrait\usebox\lycee@eqc\endgroup
  \par\smallskip}

\tikzset{grille/.style={coulGrille,line width=0.4pt}}
\newcommand{\quadrillage}[4]{\draw[grille,step=1] (#1,#2) grid (#3,#4);}
\tikzset{
  croix/.style={cross out,draw,line width=0.6pt,inner sep=0pt,minimum size=4pt},
  croixdroite/.style={croix,rotate=45,line width=0.8pt},
}
\newcommand{\pt}[3]{\path (#1) node[croix]{} node[#2,font=\small,inner sep=2.5pt]{$#3$};}

\newcommand{\lycee@repere}[4]{%
  \draw[grille,step=1] (#1,#2) grid (#3,#4);
  \draw[-{Stealth[length=2mm]},thick] (#1,0)--({#3+0.4},0) node[below left=-1pt]{$x$};
  \draw[-{Stealth[length=2mm]},thick] (0,#2)--(0,{#4+0.4}) node[below left=-1pt]{$y$};
  \node[below left,font=\scriptsize,inner sep=1.5pt] at (0,0) {$0$};}
\newcommand{\lycee@gradx}[1]{\foreach \k/\l in {#1}{%
  \draw (\k,0.07)--(\k,-0.07) node[below,font=\scriptsize,inner sep=1.5pt]{$\l$};}}
\newcommand{\lycee@grady}[1]{\foreach \k/\l in {#1}{%
  \draw (0.07,\k)--(-0.07,\k) node[left,font=\scriptsize,inner sep=1.5pt]{$\l$};}}
\newcommand{\lycee@intff}[2]{\left[#1\,;#2\right]}
\newcommand{\lycee@intfo}[2]{\left[#1\,;#2\right[}
\newcommand{\lycee@intof}[2]{\left]#1\,;#2\right]}
\newcommand{\lycee@intoo}[2]{\left]#1\,;#2\right[}
\newcommand{\lycee@ens}[1]{\left\{#1\right\}}
\AtBeginDocument{%
  \providecommand{\repere}{\lycee@repere}%
  \providecommand{\gradx}{\lycee@gradx}%
  \providecommand{\grady}{\lycee@grady}%
  \providecommand{\Cf}{\mathcal{C}\sb{f}}%
  \providecommand{\Cg}{\mathcal{C}\sb{g}}%
  \providecommand{\intff}{\lycee@intff}%
  \providecommand{\intfo}{\lycee@intfo}%
  \providecommand{\intof}{\lycee@intof}%
  \providecommand{\intoo}{\lycee@intoo}%
  \providecommand{\ens}{\lycee@ens}}

\newlist{questions}{enumerate}{3}
\setlist[questions,1]{label=\arabic*\textdegree),leftmargin=*,itemsep=2pt,topsep=3pt}
\setlist[questions,2]{label=\alph*),leftmargin=*,itemsep=1pt,topsep=2pt}
\setlist[questions,3]{label=\roman*),leftmargin=*,itemsep=1pt,topsep=1pt}
\setlist[itemize]{label=\textbullet,leftmargin=*,itemsep=2pt,topsep=3pt}

\newcommand{\besoinplace}[1]{\par\penalty\@M\begingroup
  \dimen@=#1\relax\dimen@ii=\pagegoal\advance\dimen@ii-\pagetotal
  \ifdim\dimen@>\dimen@ii\ifdim\dimen@ii>\z@\vfil\fi\break\fi\endgroup}

\newcommand{\lycee@pied}{\thepage/\pageref{LastPage}}
\newcommand{\entete}[2]{%
  \pagestyle{fancy}\fancyhf{}%
  \fancyhead[L]{\footnotesize\color{coulPart}\textbf{#1}}%
  \fancyhead[R]{\footnotesize\color{coulPart}\textbf{#2}}%
  \fancyfoot[C]{\footnotesize\lycee@pied}%
  \renewcommand{\headrulewidth}{0.4pt}%
  \renewcommand{\headrule}{\hbox to\headwidth{%
    \color{coulPart!40}\leaders\hrule height \headrulewidth\hfill}}}

\RequirePackage{ccicons}
\newcommand{\lycee@licence}{{\scriptsize\color{gray}%
  \copyright\ Luc Destombe --- \ccbyncsa\ CC BY-NC-SA 4.0}}
\newif\iflycee@licence \lycee@licencetrue
\newcommand{\sanslicence}{\lycee@licencefalse}
\AtBeginDocument{\iflycee@licence\AddToHookNext{shipout/foreground}{%
  \put(0,0){\rlap{\kern\dimexpr1in+\hoffset+\oddsidemargin+\textwidth\relax
    \llap{\raisebox{-\dimexpr1in+\voffset+\topmargin+\headheight+\headsep
      +\textheight+\footskip\relax}{\lycee@licence}}}}}\fi}

\newcommand{\lycee@site}{faire-des-maths.fr}
\newcommand{\lycee@version}{{\scriptsize\color{gray}\lycee@site
  \ifdefined\lyceeversion\ --- Version~\lyceeversion\fi}}
\AtBeginDocument{\iflycee@licence\AddToHookNext{shipout/foreground}{%
  \put(0,0){\rlap{\kern\dimexpr1in+\hoffset+\oddsidemargin\relax
    \raisebox{-\dimexpr1in+\voffset+\topmargin+\headheight+\headsep
      +\textheight+\footskip\relax}{\lycee@version}}}}\fi}

\newcommand{\formulesserrees}{%
  \setlength{\abovedisplayskip}{3pt}\setlength{\belowdisplayskip}{3pt}%
  \setlength{\abovedisplayshortskip}{2pt}\setlength{\belowdisplayshortskip}{3pt}}

\setlength{\parindent}{0pt}

\RequirePackage[implicit=false,hidelinks,pdfencoding=auto,psdextra,bookmarksopen,
  bookmarksopenlevel=1,pdfcreator={},pdfproducer={}]{hyperref}
\RequirePackage{bookmark}
\hypersetup{pdfauthor={Luc Destombe},pdfkeywords={CC BY-NC-SA 4.0}}
\newcounter{lycee@signet}
\newcommand{\signet}[3][]{\stepcounter{lycee@signet}%
  \edef\lycee@dest{\if\relax\detokenize{#1}\relax
    signet.\arabic{lycee@signet}\else#1\fi}%
  \hypertarget{\lycee@dest}{}\bookmark[level=#2,dest=\lycee@dest]{#3}}
\pdfstringdefDisableCommands{%
  \def\R{R}\def\N{N}\def\Z{Z}\def\Q{Q}\def\vect#1{#1}\def\Cf{Cf}\def\Cg{Cg}\def\fois#1{×#1}%
  \def\barre#1{#1}\def\intervalle#1#2{[#1;#2]}\def\intff#1#2{[#1;#2]}%
  \def\textsuperscript#1{#1}\def\dfrac#1#2{#1/#2}\def\frac#1#2{#1/#2}%
  \def\vec#1{#1⃗}}%
\RequirePackage[official]{eurosym}

\geometry{top=0.8cm,bottom=1.3cm,left=1cm,right=1cm,footskip=0.6cm}

\pagestyle{fancy}
\fancyhf{}
\renewcommand{\headrulewidth}{0pt}
\fancyfoot[C]{\footnotesize Page \thepage{} / \pageref{LastPage}}

\newcommand{\titrefiche}[2]{%
  \begin{center}
    \Large\textbf{#1}\\[2pt]
    \large\textbf{#2}\\[2pt]
  \end{center}
  \vspace{0.10cm}}

\setlength{\columnseprule}{0.4pt}
\setlength{\columnsep}{15pt}
\renewcommand{\columnseprulecolor}{\color{black!45}}
\setlength{\multicolsep}{2pt}

\newcounter{partiefiche}
\renewcommand{\thepartiefiche}{\Roman{partiefiche}}
\newif\ifapresbandeau
\newcommand{\lycee@bandeau}[1]{%
  \par\penalty-600
  \vspace{0.08cm}%
  \noindent\colorbox{coulPart}{%
    \parbox{\dimexpr\linewidth-2\fboxsep\relax}{%
      \color{white}\bfseries\raggedright\strut#1\strut}}%
  \nopagebreak\par\nopagebreak
  \vspace{0.06cm}\nopagebreak
  \apresbandeautrue}
\newcommand{\partiefiche}[1]{%
  \refstepcounter{partiefiche}%
  \lycee@bandeau{\signet{1}{\thepartiefiche{} --- #1}\thepartiefiche{} --- #1}}
\newcommand{\partiefichesansnum}[1]{\lycee@bandeau{\signet{1}{#1}#1}}

\newcounter{exo}
\newlength{\espaceexo}\setlength{\espaceexo}{7pt}
\newlength{\espacetitre}\setlength{\espacetitre}{2pt}
\newcommand{\exo}[1][\relax]{%
  \par
  \ifapresbandeau\apresbandeaufalse\else\penalty-150\addvspace{\espaceexo}\fi
  \refstepcounter{exo}%
  \noindent\signet[exercice-\arabic{exo}]{2}{Exercice \arabic{exo}}%
  \textbf{\color{coulExo}\underline{Exercice \theexo}}%
  \ifx\relax#1\relax\else\textbf{ : #1}\fi
  \par\nopagebreak\vspace{\espacetitre}\nopagebreak
  \ignorespaces}

\setlist[enumerate]{nosep,leftmargin=*,itemsep=2pt,topsep=2pt}
\setlist[itemize]{nosep,leftmargin=*,topsep=2pt}
\newcommand{\choix}[3]{\par\hspace*{1em}%
  \makebox[0.3\linewidth][l]{a) #1}\makebox[0.3\linewidth][l]{b) #2}c) #3}
\newcommand{\deux}[2]{\par\hspace*{0.6em}%
  \makebox[0.48\linewidth][l]{#1}#2}
\newcommand{\trois}[3]{\par\hspace*{0.6em}%
  \makebox[0.32\linewidth][l]{#1}\makebox[0.32\linewidth][l]{#2}#3}

  \renewcommand{\exo}[1][\relax]{%
    \ifapresbandeau\apresbandeaufalse\else\par\penalty-150\fi
    \refstepcounter{exo}%
    \vspace{0.05cm}%
    \noindent\signet[exercice-\arabic{exo}]{2}{Exercice \arabic{exo}}%
    \textbf{\color{coulExo}\underline{Exercice \theexo}}%
    \ifx\relax#1\relax\else\textbf{ : #1}\fi%
    \par\nopagebreak
    \ignorespaces}
  \newcommand{\rappel}[1]{%
    \par\vspace{0.06cm}\nopagebreak
    \noindent\fcolorbox{black!35}{black!5}{%
      \parbox{\dimexpr\linewidth-2\fboxsep-2\fboxrule\relax}{%
        \small\textbf{Rappel.} #1}}%
    \par\vspace{0.06cm}\nopagebreak}
  \setlist[enumerate]{nosep,leftmargin=*}
  \setlist[itemize]{nosep,leftmargin=*}
  \setlength{\multicolsep}{3pt plus 1pt minus 1pt}
  \setlength{\premulticols}{10pt}
  \setlength{\postmulticols}{10pt}
  \newlist{expr}{enumerate}{1}
  \setlist[expr]{label={\Alph*\ $=$},nosep,leftmargin=*,itemsep=0pt}

\renewenvironment{center}{\par\addvspace{3pt}\centering}{\par\addvspace{3pt}}
\formulesserrees
\AtBeginDocument{\formulesserrees}
\clubpenalty=10000
\widowpenalty=10000
\displaywidowpenalty=10000
\brokenpenalty=10000
\setlength{\parskip}{0pt}

\RequirePackage{ragged2e}
\setlength{\RaggedRightRightskip}{0pt plus 1fil}
\AtBeginDocument{\RaggedRight\binoppenalty=10000 \relpenalty=10000 }
\g@addto@macro\@parboxrestore{\RaggedRight}
\makeatother
\geometry{top=0.8cm,bottom=1.3cm,left=1cm,right=1cm,footskip=0.7cm,headheight=12pt}

\newtcblisting{python}{%
  listing only,enhanced,breakable,
  colback=black!3,colframe=black!35,boxrule=0.5pt,arc=1pt,
  left=6pt,right=6pt,top=2pt,bottom=2pt,
  listing options={language=Python,basicstyle=\ttfamily\small,
    keywordstyle=\color{coulMeth}\bfseries,columns=fullflexible,
    keepspaces=true,showstringspaces=false,aboveskip=0pt,belowskip=0pt}}
\providecommand{\Card}{\operatorname{Card}}

\begin{document}
\titrefiche{Terminale spécialité --- Dénombrement}{Fiche d'exercices du chapitre}

\begin{multicols}{2}

\partiefiche{Ensembles finis et principe additif}

\exo[Une flotte de vélos]
Une entreprise de location contrôle ses $400$ vélos électriques : $52$ ont un défaut
de freins, $120$ un défaut d'éclairage et $36$ présentent les deux défauts.
\begin{enumerate}
  \item Représenter la situation par un diagramme.
  \item Combien de vélos ne présentent aucun des deux défauts ?
  \item Combien de vélos présentent exactement un des deux défauts ?
\end{enumerate}

\exo[Course à pied et vélo]
Un club compte $90$ adhérents : $58$ pratiquent la course à pied, $41$ le vélo, et
$17$ ne pratiquent aucun des deux.
\begin{enumerate}
  \item Combien d'adhérents pratiquent au moins un des deux sports ?
  \item En déduire le nombre d'adhérents qui pratiquent les deux sports.
  \item Représenter la situation par un diagramme.
  \item Combien d'adhérents pratiquent exactement un des deux sports ?
\end{enumerate}

\exo[Multiples de 2 et de 3]
On note $E=\{1\,;2\,;\dots\,;100\}$, $A$ l'ensemble des multiples de $2$ de $E$ et
$B$ l'ensemble des multiples de $3$ de $E$.
\begin{enumerate}
  \item Calculer $\Card(A)$ et $\Card(B)$.
  \item Décrire l'ensemble $A\cap B$, puis calculer son cardinal.
  \item En déduire $\Card(A\cup B)$.
  \item Combien d'entiers de $E$ ne sont ni pairs ni multiples de $3$ ?
\end{enumerate}

\exo[Trois ateliers]
À la fête de la science, chaque visiteur interrogé a suivi des ateliers
d'astronomie, de robotique ou de photographie. On sait que :
\begin{itemize}
  \item $120$ ont suivi l'atelier d'astronomie, $100$ celui de robotique et $70$ celui
        de photographie ;
  \item $40$ ont suivi l'astronomie et la robotique, $30$ l'astronomie et la
        photographie, $25$ la robotique et la photographie ;
  \item $10$ ont suivi les trois ateliers et $15$ n'en ont suivi aucun.
\end{itemize}
\begin{enumerate}
  \item Représenter la situation par un diagramme à trois ensembles, en partant du
        centre.
  \item Combien de visiteurs ont été interrogés ?
  \item Combien ont suivi exactement un atelier ? exactement deux ?
\end{enumerate}

\partiefiche{Produit cartésien et principe multiplicatif}

\exo[Écrire un produit cartésien]
On note $E=\{a\,;b\,;c\}$ et $F=\{1\,;2\}$.
\begin{enumerate}
  \item Écrire tous les éléments de $E\times F$. Combien y en a-t-il ?
  \item Écrire tous les éléments de $F\times E$. A-t-on $E\times F=F\times E$ ?
  \item Calculer $\Card(E\times E)$ et $\Card(E\times F\times F)$.
\end{enumerate}

\exo[Plaques d'immatriculation]
Une plaque d'immatriculation comporte deux lettres, puis trois chiffres de $0$ à $9$,
puis deux lettres. Les lettres I, O et U sont exclues (confusion avec 1, 0 et V) ;
lettres et chiffres peuvent se répéter.
\begin{enumerate}
  \item Combien de lettres peut-on utiliser ?
  \item Combien de plaques différentes peut-on former ?
  \item Combien de ces plaques commencent par AB ?
\end{enumerate}

\exo[Des routes entre quatre villes]
Entre quatre villes $A$, $B$, $C$ et $D$, il y a $3$ routes de $A$ à $B$, $2$ routes de
$B$ à $C$, $4$ routes de $C$ à $D$ et $2$ routes directes de $A$ à $C$. On va toujours
vers $D$, sans revenir en arrière.
\begin{enumerate}
  \item Combien y a-t-il de trajets de $A$ à $D$ qui passent par $B$ ?
  \item Combien y a-t-il de trajets de $A$ à $D$ en tout ?
  \item Au retour, on va toujours vers $A$, sans revenir en arrière. Combien y a-t-il
        d'allers-retours de $A$ à $D$ ? et d'allers-retours dont le retour ne reprend
        pas exactement le trajet de l'aller ?
\end{enumerate}

\exo[Le menu de la cantine]
La cantine propose $4$ entrées, $5$ plats, dont $2$ végétariens, et $3$ desserts.
\begin{enumerate}
  \item Combien de menus complets (entrée, plat, dessert) peut-on composer ?
  \item Combien de ces menus ont un plat végétarien ?
  \item Une formule à deux est \og entrée et plat\fg{} ou \og plat et dessert\fg.
        Combien de formules à deux peut-on composer ?
  \item Combien de repas différents (menu complet ou formule à deux) la cantine
        propose-t-elle ?
\end{enumerate}

\partiefiche{Les $k$-uplets d'un ensemble}

\exo[Les $k$-uplets d'un petit ensemble]
On note $E=\{J\,;U\,;D\,;O\}$.
\begin{enumerate}
  \item Écrire deux triplets différents de $E$, dont l'un a une lettre répétée.
  \item Combien y a-t-il de triplets de $E$ ?
  \item Calculer $\Card(E^{6})$.
  \item Combien de $6$-uplets de $E$ commencent par $J$ ?
\end{enumerate}

\exo[Coder des caractères]
En informatique, un octet est un $8$-uplet de $\{0\,;1\}$, comme $(0\,;1\,;1\,;0\,;0\,;
1\,;0\,;1)$.
\begin{enumerate}
  \item Combien y a-t-il d'octets ? Combien d'octets comportent au moins un chiffre $1$ ?
  \item La norme ASCII code chaque caractère par un $7$-uplet de $\{0\,;1\}$.
        Peut-elle coder les $95$ caractères de l'anglais ?
  \item Les langues européennes demandent $191$ caractères. Quel est le plus petit
        entier $k$ tel que des $k$-uplets de $\{0\,;1\}$ suffisent ?
\end{enumerate}

\exo[Le code génétique]
Un codon est une suite de trois nucléotides, chacun choisi parmi A, C, G et T ; l'ordre
compte. Chaque codon code l'un des $20$ acides aminés, ou la fin d'une protéine.
\begin{enumerate}
  \item Combien y a-t-il de codons ? Y en a-t-il assez pour les $20$ acides aminés ?
  \item Des suites de deux nucléotides auraient-elles suffi ?
  \item Combien de codons ne comportent pas de G ?
  \item En déduire le nombre de codons qui comportent au moins un G.
\end{enumerate}

\exo[Codes de cinq chiffres]
Un code est formé de $5$ chiffres de $0$ à $9$, distincts ou non (par exemple $0\,4\,4\,7\,0$).
\begin{enumerate}
  \item Combien y a-t-il de codes ?
  \item Combien de codes se terminent par $0$ ou par $5$ ?
  \item Combien de codes ne comportent aucun chiffre $7$ ?
  \item En déduire le nombre de codes qui comportent au moins un $7$.
\end{enumerate}

\exo[Colorier un quadrillage]
On colorie en rouge, jaune ou noir chacune des $16$ cases d'un quadrillage.
\begin{enumerate}
  \item Combien y a-t-il de coloriages ?
  \item Combien de coloriages n'utilisent pas le noir ?
  \item Combien de coloriages comportent au moins une case noire ?
  \item Combien de coloriages comportent exactement une case noire ?
\end{enumerate}

\exo[La force d'un mot de passe]
Un mot de passe est formé de $4$ caractères, chacun choisi parmi les $26$ lettres
minuscules et les $10$ chiffres.
\begin{enumerate}
  \item Combien y a-t-il de mots de passe ?
  \item Combien de mots de passe ne comportent que des lettres ? au moins un chiffre ?
  \item Un logiciel essaie $1\,000$ mots de passe par seconde. Combien de temps lui
        faut-il, au plus, pour les essayer tous ?
  \item On passe à $8$ caractères, choisis parmi $62$ (minuscules, majuscules,
        chiffres). Combien de jours faut-il, au plus, à un logiciel qui essaie un
        milliard de mots de passe par seconde ?
\end{enumerate}

\partiefiche{Les $k$-uplets d'éléments distincts, factorielle}

\exo[Calculer avec des factorielles]
\begin{enumerate}
  \item Calculer $6!$ et $8!$.
  \item Calculer sans calculatrice :
        \begin{center}
        $A=\dfrac{10!}{8!}$ \qquad $B=\dfrac{12!}{9!}$ \qquad $C=\dfrac{7!}{3!\times 4!}$
        \end{center}
  \item Soit $n\in\N$ tel que $n\geq 2$. Simplifier :
        \begin{center}
        $D=\dfrac{(n+1)!}{n!}$ \qquad $F=\dfrac{n!}{(n-2)!}$
        \end{center}
\end{enumerate}

\exo[Triplets d'éléments distincts]
On note $E=\{a\,;b\,;c\,;d\,;e\}$.
\begin{enumerate}
  \item Écrire trois triplets d'éléments distincts de $E$.
  \item Combien y a-t-il de triplets d'éléments distincts de $E$ ?
  \item Combien de ces triplets commencent par $a$ ?
  \item Combien y a-t-il de couples d'éléments distincts de $E$ ?
\end{enumerate}

\exo[Bac 2025, Amérique du Sud]
Une association organise une compétition de course de haies qui permettra d'établir
un podium (le podium est constitué des trois meilleurs sportifs classés dans leur
ordre d'arrivée). Sept sportifs participent au tournoi. Jacques est l'un d'entre eux.
\begin{enumerate}
  \item Combien de podiums différents sont possibles ?
  \item \textbf{Affirmation :} il y a $90$ podiums différents dont Jacques fait
        partie. Dire si l'affirmation est vraie ou fausse, en justifiant.
\end{enumerate}

\exo[Bac 2025, Métropole]
Lorsqu'une personne achète son billet de musée en ligne, un code de validation lui est
envoyé par SMS afin qu'elle confirme son achat. Ce code est généré de façon aléatoire
et est constitué de $4$ chiffres deux à deux distincts, le premier chiffre étant
différent de $0$.

\textbf{Affirmation :} le nombre de codes différents pouvant être générés est
$5\,040$. Dire si l'affirmation est vraie ou fausse, en justifiant.

\exo[La soirée théâtre]
La soirée théâtre du lycée propose $8$ saynètes différentes ; $5$ sont jouées avant
l'entracte, les $3$ autres après.
\begin{enumerate}
  \item On tire au sort les $5$ saynètes d'avant l'entracte et leur ordre. Combien de
        premières parties sont possibles ?
  \item Combien de ces premières parties commencent par la saynète de la classe de
        terminale ?
  \item Finalement, on sait quelles saynètes sont jouées avant l'entracte, mais pas
        encore dans quel ordre, ni l'ordre de celles d'après. Combien de soirées
        sont encore possibles ?
\end{enumerate}

\exo[Codes de trois lettres]
Un code est formé de $3$ lettres de l'alphabet ($26$ lettres).
\begin{enumerate}
  \item Combien y a-t-il de codes ?
  \item Combien de codes ont leurs $3$ lettres différentes ?
  \item En déduire le nombre de codes qui ont au moins deux lettres identiques.
\end{enumerate}

\partiefiche{Les permutations}

\exo[Anagrammes]
Une anagramme d'un mot est un mot (qui a un sens ou non) écrit avec les mêmes lettres,
chacune autant de fois.
\begin{enumerate}
  \item Combien le mot JUDO a-t-il d'anagrammes ? et le mot BRAVO ?
  \item Combien d'anagrammes de BRAVO commencent par une voyelle ?
  \item \textbf{Affirmation (bac 2025, Nouvelle-Calédonie) :} le mot EULER possède
        $120$ anagrammes. Dire si l'affirmation est vraie ou fausse, en justifiant.
\end{enumerate}

\exo[Des lettres qui se répètent]
\begin{enumerate}
  \item On écrit les $5$ lettres de KAYAK dans tous les ordres possibles. Expliquer
        pourquoi chaque anagramme est obtenue $2!\times 2!$ fois, puis compter les
        anagrammes de KAYAK.
  \item Compter de même les anagrammes de BANANE.
  \item \textbf{Bac 2026, Nouvelle-Calédonie.} On s'intéresse au nombre $N$ de mots
        de six lettres (ayant un sens ou non) que l'on peut écrire avec les lettres A,
        B, C, D, E, en utilisant exactement deux fois la lettre A, et exactement une
        fois chacune des autres lettres. Exemples : AABCDE ; BACADE ; CDBAEA.

        \textbf{Affirmation :} $N=120$. Dire si l'affirmation est vraie ou fausse, en
        justifiant.
\end{enumerate}

\exo[La photo de famille]
Une famille de six personnes, deux parents et quatre enfants, s'aligne pour une photo.
\begin{enumerate}
  \item Combien y a-t-il d'alignements possibles ?
  \item Combien d'alignements placent les parents aux deux extrémités ?
  \item Combien d'alignements placent les parents côte à côte ?
  \item En déduire le nombre d'alignements où les parents ne sont pas côte à côte.
\end{enumerate}

\exo[Factorielle en Python]
On considère la fonction \texttt{mystere} définie ci-dessous en Python :
\begin{python}
def mystere(n):
    f = 1
    for k in range(1, n+1):
        f = f * k
    return f
\end{python}
\begin{enumerate}
  \item \textbf{Affirmation (bac 2026, Polynésie) :} le nombre renvoyé lors de
        l'exécution de \texttt{mystere(5)} est $120$. Vrai ou faux ? Justifier.
  \item Que renvoie \texttt{mystere(n)} pour $n\in\Ns$ ?
\end{enumerate}

\partiefiche{Les combinaisons}

\exo[Calculer et écrire des combinaisons]
\begin{enumerate}
  \item Calculer sans calculatrice :
        \begin{center}
        $\dbinom{5}{2}$ \qquad $\dbinom{7}{3}$ \qquad $\dbinom{10}{1}$ \qquad
        $\dbinom{12}{0}$ \qquad $\dbinom{20}{3}$
        \end{center}
  \item On note $E=\{a\,;b\,;c\,;d\,;e\}$. Écrire toutes les parties de $E$ à $2$
        éléments, et vérifier qu'il y en a $\dbinom{5}{2}$.
  \item Combien $E$ a-t-il de parties à $3$ éléments ?
\end{enumerate}

\exo[Bac 2025, Amérique du Sud]
Dire si chaque affirmation est vraie ou fausse, en justifiant.
\begin{enumerate}
  \item Deux équipes de footballeurs de $22$ et $25$ joueurs échangent une poignée de
        main à la fin d'un match. Chaque joueur d'une équipe serre une seule fois la
        main de chaque joueur de l'autre équipe.

        \textbf{Affirmation 1 :} $47$ poignées de mains ont été échangées.
  \item Une course oppose $18$ concurrents. On récompense indistinctement les trois
        premiers en offrant le même prix à chacun.

        \textbf{Affirmation 2 :} il y a $4\,896$ possibilités de distribuer ces prix.
  \item Dans un groupe de $12$ personnes, chacune serre une fois la main de chacune des
        autres.

        \textbf{Affirmation 3 :} $66$ poignées de mains ont été échangées.
\end{enumerate}

\exo[Bac 2025 et 2026, former des équipes]
Dire si chaque affirmation est vraie ou fausse, en justifiant.
\begin{enumerate}
  \item \emph{(Polynésie 2025)} Dans une classe de terminale, il y a $18$ filles et
        $14$ garçons. On constitue une équipe de volley-ball en choisissant au hasard
        $3$ filles et $3$ garçons.

        \textbf{Affirmation 1 :} il y a $297\,024$ possibilités pour former une telle
        équipe.
  \item \emph{(Métropole 2026)} Un jeu sur le thème de la culture de la Pop est proposé
        à $10$ personnes dont $4$ sont des musiciens professionnels. Deux équipes de $5$
        joueurs s'affrontent. Elles sont chacune constituées de $2$ musiciens
        professionnels et de $3$ personnes qui ne le sont pas.

        \textbf{Affirmation 2 :} on peut former ainsi $120$ équipes différentes.
\end{enumerate}

\exo[Bac 2026, Asie]
Une urne contient $32$ jetons numérotés de $1$ à $32$ indiscernables au toucher. On
tire simultanément $5$ jetons de cette urne. On appelle tirage la liste non ordonnée
des numéros des cinq jetons tirés.
\begin{enumerate}
  \item Combien y a-t-il de tirages possibles ?
  \item \textbf{Affirmation :} le nombre de tirages possibles contenant au moins un
        multiple de $8$ est égal à $103\,096$. Vrai ou faux ? Justifier.
\end{enumerate}

\exo[Le club de robotique]
Le club de robotique compte $16$ membres : $9$ élèves de première et $7$ de terminale.
Il envoie une délégation de $5$ membres à une compétition.
\begin{enumerate}
  \item Combien de délégations peut-on former ?
  \item Combien de délégations comptent $3$ élèves de première et $2$ de terminale ?
  \item Combien de délégations comptent au moins un élève de terminale ?
  \item La délégation choisie, on attribue à ses $5$ membres $5$ rôles différents
        (pilote, codeur\dots). Combien d'attributions sont possibles ?
  \item Combien y a-t-il de délégations avec leurs rôles ? Retrouver ce nombre à l'aide
        des questions 1 et 4.
\end{enumerate}

\exo[Bac 2026, Polynésie]
Un refuge spécialisé accueille $24$ rapaces numérotés de $1$ à $24$. Le refuge
travaille avec trois vétérinaires. Afin d'administrer un soin à tous les oiseaux, le
directeur veut donc former trois groupes, numérotés de $1$ à $3$, constitués chacun
de $8$ rapaces.

\textbf{Affirmation :} le nombre de façons de répartir les rapaces dans les groupes
$1$ à $3$ est $\dfrac{24!}{(8!)^{3}}$. Dire si l'affirmation est vraie ou fausse, en
justifiant.

\partiefiche{Symétrie, relation et triangle de Pascal}

\exo[Prolonger le triangle de Pascal]
\begin{enumerate}
  \item Écrire le triangle de Pascal jusqu'à la ligne $n=8$.
  \item Y lire $\dbinom{8}{3}$ et $\dbinom{8}{4}$.
  \item Calculer la somme des nombres de la ligne $n=8$. Quel nombre reconnaît-on ?
\end{enumerate}

\exo[Symétrie]
\begin{enumerate}
  \item Calculer sans calculatrice :
        \begin{center}
        $\dbinom{50}{49}$ \qquad $\dbinom{100}{98}$ \qquad $\dbinom{15}{13}$
        \end{center}
  \item Déterminer l'entier $n\geq 2$ tel que $\dbinom{n}{2}=45$.
\end{enumerate}

\exo[Démontrer des égalités]
\begin{enumerate}
  \item \textbf{Bac 2025, Polynésie.} Soient $n$ et $k$ deux entiers naturels non nuls
        tels que $k\leq n$.

        \textbf{Affirmation :} $n\times\dbinom{n-1}{k-1}=k\times\dbinom{n}{k}$. Vrai ou
        faux ? Justifier.
  \item Soit $n\in\N$ tel que $n\geq 2$.
        \begin{enumerate}[label=\alph*)]
          \item Vérifier que $\dbinom{n}{2}+\dbinom{n+1}{2}=n^{2}$ pour $n=3$ et $n=4$.
          \item Le démontrer pour tout $n\geq 2$.
        \end{enumerate}
\end{enumerate}

\partiefiche{Nombre de parties d'un ensemble, mots et chemins}

\exo[Toutes les parties]
On note $E=\{a\,;b\,;c\,;d\}$.
\begin{enumerate}
  \item Combien $E$ a-t-il de parties ?
  \item Écrire toutes ces parties, rangées selon leur nombre d'éléments.
  \item Vérifier que $\dbinom{4}{0}+\dbinom{4}{1}+\dbinom{4}{2}+\dbinom{4}{3}+
        \dbinom{4}{4}=2^{4}$.
\end{enumerate}

\exo[Parties et $n$-uplets de $\{0\,;1\}$]
On note $E=\{1\,;2\,;3\,;4\,;5\,;6\}$. À chaque partie de $E$, on associe le $6$-uplet
de $\{0\,;1\}$ dont le chiffre numéro $i$ vaut $1$ si $i$ est dans la partie, et $0$
sinon.
\begin{enumerate}
  \item Combien $E$ a-t-il de parties ?
  \item Quel $6$-uplet correspond à la partie $\{1\,;3\,;6\}$ ? et à $\varnothing$ ?
  \item Quelle partie correspond au $6$-uplet $(0\,;1\,;1\,;1\,;0\,;0)$ ?
  \item Combien de $6$-uplets de $\{0\,;1\}$ comportent exactement deux chiffres $1$ ?
\end{enumerate}

\exo[L'écriture braille]
Un caractère braille est une grille de $6$ points, chacun en relief ou non. La grille
sans aucun point en relief représente un espace, pas un caractère.
\begin{enumerate}
  \item Combien de caractères braille peut-on former ?
  \item Combien de caractères comportent exactement $3$ points en relief ?
  \item Combien de caractères comportent au moins $2$ points en relief ?
\end{enumerate}

\exo[Lancers d'une pièce]
On lance $6$ fois une pièce et on note la suite des résultats, P (pile) ou F (face).
\begin{enumerate}
  \item Combien y a-t-il de suites de résultats ? Combien de chemins compte l'arbre
        qui les représente ?
  \item Combien de suites comportent exactement $2$ fois pile ?
  \item Combien de suites comportent au moins une fois pile ?
  \item Combien de suites comportent autant de piles que de faces ?
\end{enumerate}

\exo[Les chemins d'un robot]
\begin{minipage}[t]{0.55\linewidth}\raggedright
Un robot part du point $O$ et rejoint le point $B$ en suivant les lignes du
quadrillage, à chaque pas d'une case vers la droite (D) ou vers le haut (H).
\end{minipage}\hfill
\begin{minipage}[t]{0.42\linewidth}\centering\vspace{-2pt}
\begin{tikzpicture}[x=0.5cm,y=0.5cm]
  \draw[grille] (0,0) grid[step=1] (5,3);
  \path (0,0) node[croix]{} node[below left,font=\footnotesize]{$O$}
        (5,3) node[croix]{} node[above right,font=\footnotesize]{$B$}
        (2,1) node[croix]{} node[above left,font=\footnotesize]{$C$};
\end{tikzpicture}
\end{minipage}
\begin{enumerate}
  \item Expliquer pourquoi un chemin de $O$ à $B$ est un mot de $8$ lettres qui
        comporte $5$ fois D et $3$ fois H.
  \item Combien y a-t-il de chemins de $O$ à $B$ ?
  \item Combien de ces chemins passent par le point $C$ ?
  \item Le robot choisit un chemin au hasard. Quelle est la probabilité qu'il passe
        par $C$ ?
\end{enumerate}

\partiefiche{Choisir le bon modèle}

\exo[Quel objet compter ?]
Pour chaque situation, dire quel objet on compte ($k$-uplet, $k$-uplet d'éléments
distincts, permutation, combinaison), puis faire le calcul.
\begin{enumerate}[label=\alph*)]
  \item Le classement des trois premiers d'une course de $20$ cyclistes.
  \item Le choix de $5$ livres parmi $12$ pour partir en vacances.
  \item Un code de $4$ lettres de l'alphabet.
  \item L'ordre de passage de $8$ élèves à l'oral.
  \item L'attribution de $3$ prix différents à $10$ élèves, un élève pouvant
        recevoir plusieurs prix.
  \item Un groupe de $3$ élèves choisis parmi $30$.
\end{enumerate}

\exo[QCM]
Pour chaque question, une seule réponse est exacte ; la donner en justifiant.
\begin{enumerate}
  \item \emph{(Bac 2026, Amérique du Sud)} Une école de rugby compte $20$ joueurs. Les
        matchs auxquels peut participer cette école se jouent à $10$. Combien d'équipes
        différentes cette école peut-elle former ?

        a)~$\dbinom{20}{10}$ \quad b)~$20^{10}$ \quad
        c)~$20\times 19\times 18\times\dots\times 11$ \quad d)~$10^{20}$
  \item Le nombre d'anagrammes du mot NOTE est :

        a)~$4^{4}$ \quad b)~$24$ \quad c)~$4$ \quad d)~$\dbinom{4}{4}$
  \item Le nombre de parties d'un ensemble à $5$ éléments est :

        a)~$25$ \quad b)~$10$ \quad c)~$32$ \quad d)~$120$
  \item Le nombre de façons d'asseoir $4$ personnes sur $6$ chaises alignées est :

        a)~$\dbinom{6}{4}$ \quad b)~$6^{4}$ \quad c)~$4!$ \quad d)~$360$
\end{enumerate}

\exo[Trois dés]
On lance trois dés équilibrés à six faces, de couleurs différentes (rouge, vert, bleu),
et on note le triplet des résultats.
\begin{enumerate}
  \item Combien y a-t-il d'issues ?
  \item Calculer la probabilité que les trois numéros soient différents.
  \item Calculer la probabilité d'obtenir au moins un $6$.
  \item Calculer la probabilité que la somme des trois numéros soit égale à $4$.
\end{enumerate}

\exo[Une poignée de jetons]
Un sac contient $12$ jetons : $5$ rouges, $4$ bleus et $3$ verts. On en tire $4$
simultanément, au hasard.
\begin{enumerate}
  \item Combien y a-t-il de tirages ?
  \item Calculer la probabilité que les $4$ jetons soient de la même couleur.
  \item Calculer la probabilité d'obtenir exactement $2$ jetons rouges.
  \item Calculer la probabilité d'obtenir au moins un jeton vert.
\end{enumerate}

\partiefichesansnum{Exercices type bac}

\exo[Bac 2026, Nouvelle-Calédonie]
Un village organise une fête.

\textbf{Partie A}

Les organisateurs de la fête recherchent $3$ personnes pour tenir la buvette. $15$
villageois se portent volontaires pour cette mission.
\begin{enumerate}
  \item Combien existe-t-il de manières de choisir les $3$ personnes qui tiendront la
        buvette parmi les $15$ villageois volontaires ?
  \item Parmi les $15$ villageois volontaires, il y a $5$ hommes et $10$ femmes. Les
        organisateurs souhaitent à présent que parmi les $3$ personnes qui tiendront la
        buvette, il y ait $1$ homme et $2$ femmes. En respectant le souhait des
        organisateurs, combien existe-t-il de manières de choisir les $3$ personnes qui
        tiendront la buvette parmi les $15$ villageois volontaires ?
\end{enumerate}

\textbf{Partie B}

Les participants à la fête peuvent habiter le village ou non, consommer une boisson ou
non. On dispose des informations suivantes :
\begin{itemize}
  \item $20\,\%$ des participants à la fête habitent le village. Parmi eux, les trois
        quarts ont consommé une boisson ;
  \item la moitié des participants ont consommé une boisson.
\end{itemize}
On choisit au hasard un participant à la fête et on note les événements suivants :
$V$ : \og le participant habite le village\fg{} ; $B$ : \og le participant consomme
une boisson\fg.
\begin{enumerate}
  \item L'arbre pondéré ci-dessous représente la situation. Recopier cet arbre et
        compléter les pointillés.
        \begin{center}
        \begin{tikzpicture}[x=1.6cm,y=0.55cm,font=\small]
          \draw (0,0) -- node[above left]{$\dots$} (1,1.5) node[right]{$V$};
          \draw (0,0) -- node[below left]{$\dots$} (1,-1.5) node[right]{$\overline{V}$};
          \draw (1.35,1.5) -- node[above]{$\dots$} (2.4,2.4) node[right]{$B$};
          \draw (1.35,1.5) -- node[below]{$\dots$} (2.4,0.6) node[right]{$\overline{B}$};
          \draw (1.35,-1.5) -- (2.4,-0.6) node[right]{$B$};
          \draw (1.35,-1.5) -- (2.4,-2.4) node[right]{$\overline{B}$};
        \end{tikzpicture}
        \end{center}
  \item Déterminer la probabilité que le participant habite le village et consomme une
        boisson.
  \item On note $P_{\overline{V}}(B)$ la probabilité que le participant consomme une
        boisson sachant qu'il n'habite pas le village. Montrer que
        $P_{\overline{V}}(B)=\dfrac{7}{16}$.
\end{enumerate}

\exo[Bac 2026, Amérique du Sud]
Un sac opaque contient dix boules indiscernables au toucher et numérotées de $0$ à $9$.
Une partie d'un jeu consiste à tirer deux boules dans ce sac, successivement et sans
remise, pour former un nombre à deux chiffres dont le premier chiffre est le numéro
de la première boule tirée et le deuxième chiffre est le numéro de la deuxième boule
tirée. Par exemple si l'on tire successivement la boule $6$ puis la boule $2$ le nombre
formé est $62$.
\begin{enumerate}
  \item Justifier que l'on peut former $90$ nombres différents à l'aide de ce processus.
  \item Pour gagner à ce jeu il faut réaliser l'événement $S$ : \og le nombre à deux
        chiffres ainsi formé contient deux chiffres tous deux inférieurs ou égaux à
        $5$\fg.
        \begin{enumerate}[label=\alph*)]
          \item Combien de nombres gagnants sont possibles à l'issue de ce jeu ?
                Justifier.
          \item On note $P(S)$ la probabilité de l'événement $S$. Montrer que
                $P(S)=\dfrac{1}{3}$.
        \end{enumerate}
  \item Parmi tous les nombres à deux chiffres que l'on peut former en jouant à ce jeu,
        combien ont au moins un des deux chiffres le composant qui soit supérieur ou
        égal à $8$ ?
\end{enumerate}

\exo[Bac 2025, Centres étrangers]
Le codage \og base64\fg, utilisé en informatique, permet de représenter et de
transmettre des messages et d'autres données telles que des images, en utilisant $64$
caractères : les $26$ lettres majuscules, les $26$ lettres minuscules, les chiffres de
$0$ à $9$ et deux autres caractères spéciaux.

On s'intéresse aux séquences de $4$ caractères en base64. Par exemple, \og gP3g\fg{}
est une telle séquence. Dans une séquence, l'ordre est à prendre en compte : les
séquences \og m5C2\fg{} et \og 5C2m\fg{} ne sont pas identiques.
\begin{enumerate}
  \item Déterminer le nombre de séquences possibles.
  \item Déterminer le nombre de séquences si l'on impose que les $4$ caractères sont
        différents deux à deux.
  \item
        \begin{enumerate}[label=\alph*)]
          \item Déterminer le nombre de séquences ne comportant pas de lettre A
                majuscule.
          \item En déduire le nombre de séquences comportant au moins une lettre A
                majuscule.
          \item Déterminer le nombre de séquences comportant exactement une fois la
                lettre A majuscule.
          \item Déterminer le nombre de séquences comportant exactement deux fois la
                lettre A majuscule.
        \end{enumerate}
\end{enumerate}

\end{multicols}

\end{document}
